\chapter*{Acronyms and Key Abbreviations}
\addcontentsline{toc}{chapter}{Acronyms and Key Abbreviations}

This section consolidates the main acronyms used throughout the book. The goal is not merely to expand abbreviations, but to remind the reader of the system or managerial meaning of each term in the e-commerce context.

\begin{description}
  \item[AI] Artificial Intelligence. The broader field concerned with building systems that perform tasks associated with human intelligence.
  \item[AML] Anti-Money Laundering. Controls and monitoring processes designed to detect and prevent suspicious financial activity.
  \item[API] Application Programming Interface. A contract through which one software component exposes services or data to another.
  \item[AUC] Area Under the Curve. A performance metric often used in binary classification, especially fraud and risk models.
  \item[B2B] Business-to-Business. Commerce transactions where the buyer is another business entity.
  \item[B2C] Business-to-Consumer. Commerce transactions where the seller serves end consumers directly.
  \item[BI] Business Intelligence. Reporting and analytical practices for monitoring business performance and supporting decisions.
  \item[CI/CD] Continuous Integration and Continuous Delivery or Deployment. Engineering practices for automating build, test, and release workflows.
  \item[CLV] Customer Lifetime Value. The expected long-term commercial value generated by a customer relationship.
  \item[CRM] Customer Relationship Management. Systems and processes for managing customer data, interactions, and service workflows.
  \item[D2C] Direct-to-Consumer. A business model where brands sell directly to end customers without relying solely on intermediaries.
  \item[DL] Deep Learning. A subset of machine learning based on multi-layer neural network models.
  \item[EA] Enterprise Architecture. The discipline of aligning business capabilities, processes, systems, and technology strategy.
  \item[ERP] Enterprise Resource Planning. Integrated enterprise software used to manage operational processes such as finance, inventory, procurement, and order management.
  \item[ESB] Enterprise Service Bus. An integration approach centered on mediation, routing, transformation, and orchestration between systems.
  \item[ETL] Extract, Transform, Load. A data integration pattern where data is moved and transformed before being stored in a target system.
  \item[ELT] Extract, Load, Transform. A data integration pattern where raw data is loaded first and transformed inside the target platform.
  \item[GDPR] General Data Protection Regulation. A European privacy regulation governing personal data processing and data subject rights.
  \item[HTTP] Hypertext Transfer Protocol. The application protocol underlying most web communication.
  \item[ID] Identifier. A key or token used to distinguish a business entity, user, product, order, or session.
  \item[iPaaS] Integration Platform as a Service. A managed integration environment used to connect applications, data flows, and enterprise processes.
  \item[KPI] Key Performance Indicator. A measurable indicator used to assess business or operational performance.
  \item[LLM] Large Language Model. A generative AI model trained on large-scale language data and used for tasks such as summarization, dialogue, and content generation.
  \item[LTR] Learning to Rank. Machine learning methods designed to optimize ranked outputs such as search or recommendation lists.
  \item[MAE] Mean Absolute Error. A common regression metric for forecasting or pricing models.
  \item[MCQ] Multiple-Choice Question. A structured assessment format used throughout this book.
  \item[ML] Machine Learning. Methods that learn statistical patterns from data to support prediction, classification, ranking, or decision assistance.
  \item[MLOps] Machine Learning Operations. Practices for deploying, monitoring, governing, and maintaining production ML systems.
  \item[NDCG] Normalized Discounted Cumulative Gain. A ranking metric commonly used in recommendation and search evaluation.
  \item[NLP] Natural Language Processing. Methods for analyzing, understanding, and generating human language.
  \item[ORM] Operations Research and Management Science. Quantitative methods for optimization, planning, and decision support.
  \item[PDF] Portable Document Format. The format used for the compiled output of this manuscript.
  \item[PIM] Product Information Management. Processes and systems for managing product attributes, descriptions, taxonomy, and enrichment.
  \item[POS] Point of Sale. Physical or digital checkout systems that record transactions and payments.
  \item[RAG] Retrieval-Augmented Generation. An LLM architecture pattern in which external documents are retrieved and injected into model context.
  \item[RBAC] Role-Based Access Control. An access model where permissions are assigned according to defined organizational roles.
  \item[RMSE] Root Mean Squared Error. A regression metric often used in forecasting and pricing evaluation.
  \item[ROI] Return on Investment. A business measure of value relative to cost.
  \item[SCM] Supply Chain Management. Processes and systems for planning, sourcing, fulfillment, and movement of goods across the supply network.
  \item[SLA] Service Level Agreement. A defined commitment regarding service performance, availability, or response time.
  \item[SLO] Service Level Objective. A target service quality threshold used in operational engineering.
  \item[SKU] Stock Keeping Unit. A distinct inventory item identifier used for catalog and inventory control.
  \item[SQL] Structured Query Language. The dominant language for querying relational data systems.
  \item[SSO] Single Sign-On. An authentication pattern that allows users to access multiple systems through one identity flow.
  \item[TCO] Total Cost of Ownership. The full lifecycle cost of acquiring, operating, integrating, and maintaining a solution.
  \item[TOC] Table of Contents. The navigational chapter list in the front matter of the book.
  \item[UI] User Interface. The visible interaction layer through which a user engages with a system.
  \item[UX] User Experience. The broader quality of a user’s interaction with a system, including usability, trust, speed, and clarity.
  \item[WMS] Warehouse Management System. Software and processes used to manage storage, picking, packing, movement, and dispatch of inventory.
\end{description}

\section*{How to Use This Section}
Readers should treat this glossary as a practical reference, especially when moving between technical and managerial chapters.
The same acronym may appear in different operational contexts across the book, but its meaning here provides the baseline interpretation intended by the manuscript.
